% p156_1; source PDF page 156; manually verified editable table.
\begin{tikzpicture}[x=1cm,y=-1cm,inner sep=0pt,outer sep=0pt]
\path[use as bounding box] (0,0) rectangle (16.000,9.780);
\draw[line width=.4pt] (0.000,0.000) rectangle (2.300,0.855);
\node[align=center,text width=2.040cm,font={\linespread{1}\fontsize{10}{13.50}\selectfont\bfseries}] at (1.150,0.427) {参与时间};
\draw[line width=.4pt] (2.300,0.000) rectangle (6.500,0.855);
\node[align=center,text width=3.940cm,font={\linespread{1}\fontsize{10}{13.50}\selectfont\bfseries}] at (4.400,0.427) {结论};
\draw[line width=.4pt] (6.500,0.000) rectangle (16.000,0.855);
\node[align=center,text width=9.240cm,font={\linespread{1}\fontsize{10}{13.50}\selectfont\bfseries}] at (11.250,0.427) {范例};
\draw[line width=.4pt] (0.000,0.855) rectangle (2.300,2.460);
\node[align=center,text width=2.040cm,font={\linespread{1}\fontsize{10}{13.50}\selectfont}] at (1.150,1.657) {事前共谋，临时迟到};
\draw[line width=.4pt] (2.300,0.855) rectangle (6.500,2.460);
\node[align=center,text width=3.940cm,font={\linespread{1}\fontsize{10}{13.50}\selectfont}] at (4.400,1.657) {成立共同犯罪};
\draw[line width=.4pt] (6.500,0.855) rectangle (16.000,2.460);
\node[align=left,text width=9.240cm,font={\linespread{1}\fontsize{10}{13.50}\selectfont}] at (11.250,1.657) {甲乙共谋明晚抢劫丙，甲先到，重伤丙，乙赶到，拿走丙财物。甲乙构成抢劫罪（致人重伤）的共同犯罪（2018 年试题）};
\draw[line width=.4pt] (0.000,2.460) rectangle (2.300,4.965);
\node[align=center,text width=2.040cm,font={\linespread{1}\fontsize{10}{13.50}\selectfont}] at (1.150,3.712) {中途参与};
\draw[line width=.4pt] (2.300,2.460) rectangle (6.500,4.965);
\node[align=center,text width=3.940cm,font={\linespread{1}\fontsize{10}{13.50}\selectfont}] at (4.400,3.712) {成立承继的共同犯罪};
\draw[line width=.4pt] (6.500,2.460) rectangle (16.000,4.965);
\node[align=left,text width=9.240cm,font={\linespread{1}\fontsize{10}{13.50}\selectfont}] at (11.250,3.712) {甲抢劫丙，重伤丙后，乙才参与，拿走丙的财物，交给甲。甲构成抢劫罪（致人重伤）。乙构成甲的抢劫罪的承继共同犯罪，但对重伤不负责，因为重伤是乙参与之前甲制造的（2007 年第 53 题）②};
\draw[line width=.4pt] (0.000,4.965) rectangle (2.300,9.780);
\node[align=center,text width=2.040cm,font={\linespread{1}\fontsize{10}{13.50}\selectfont}] at (1.150,7.373) {事后参与（事后帮）③};
\draw[line width=.4pt] (2.300,4.965) rectangle (6.500,6.570);
\node[align=left,text width=3.940cm,font={\linespread{1}\fontsize{10}{13.50}\selectfont}] at (4.400,5.768) {帮助对象是犯罪人，构成窝藏、包庇罪};
\draw[line width=.4pt] (6.500,4.965) rectangle (16.000,6.570);
\node[align=left,text width=9.240cm,font={\linespread{1}\fontsize{10}{13.50}\selectfont}] at (11.250,5.768) {甲抢劫到丙的财物后，乙骑车路过，载甲离去。乙构成窝藏罪};
\draw[line width=.4pt] (2.300,6.570) rectangle (6.500,8.175);
\node[align=left,text width=3.940cm,font={\linespread{1}\fontsize{10}{13.50}\selectfont}] at (4.400,7.373) {帮助对象是赃物，构成掩饰、隐瞒犯罪所得罪};
\draw[line width=.4pt] (6.500,6.570) rectangle (16.000,8.175);
\node[align=left,text width=9.240cm,font={\linespread{1}\fontsize{10}{13.50}\selectfont}] at (11.250,7.373) {甲抢劫到丙的财物后，乙帮助销赃。乙构成掩饰、隐瞒犯罪所得罪};
\draw[line width=.4pt] (2.300,8.175) rectangle (6.500,9.780);
\node[align=left,text width=3.940cm,font={\linespread{1}\fontsize{10}{13.50}\selectfont}] at (4.400,8.978) {帮助对象是证据，构成帮助毁灭、伪造证据罪};
\draw[line width=.4pt] (6.500,8.175) rectangle (16.000,9.780);
\node[align=left,text width=9.240cm,font={\linespread{1}\fontsize{10}{13.50}\selectfont}] at (11.250,8.978) {甲抢劫到丙的财物后，乙帮助销毁抢劫用的工具。乙构成帮助毁灭证据罪};
\end{tikzpicture}
